% =============================================================================
%  CHAPTER 1 — INTRODUCTION (~5% of thesis, ~2-3 pages)
% =============================================================================
\section{Introduction}
\label{sec:introduction}

\subsection{Motivation}
\label{subsec:motivation}

The Tunisian Baccalaureate (Baccalaur\'{e}at) mathematics examination,
particularly for the Section Math\'{e}matiques track, is a high-stakes
national assessment that determines university admission for thousands of
students each year. Unlike standardized tests in other countries, the Tunisian
Bac demands not only correct mathematical reasoning but also strict adherence
to a specific \textit{correction redaction style}—a formal writing convention
that dictates how proofs, derivations, and solutions must be structured and
phrased in French. Students who produce correct answers but fail to follow
this style risk losing significant marks.

Despite the widespread adoption of large language models (LLMs) in
educational technology~\cite{kasneci2023chatgpt_education}, no existing AI
tutoring system addresses the specific needs of Tunisian Bac mathematics
students. General-purpose chatbots such as ChatGPT lack knowledge of the
Tunisian curriculum boundaries—they may suggest methods from university-level
mathematics that are explicitly outside the Bac program—and they cannot
reproduce the distinctive correction style expected by Tunisian examiners.

This gap motivates the central question of this thesis: \textbf{Can a
Retrieval-Augmented Generation system, grounded in a curated corpus of
Tunisian Bac materials, produce curriculum-faithful answers that respect the
official correction style?}

\subsection{Problem Statement}
\label{subsec:problem-statement}

Building an effective AI tutor for this domain requires addressing three
interrelated challenges:

\begin{enumerate}
    \item \textbf{Curriculum fidelity.} The system must answer exclusively
    within the boundaries of the Tunisian Bac program. Using methods from
    outside the curriculum (e.g., university-level analysis or algebra) is
    not only unhelpful but actively harmful to exam preparation.

    \item \textbf{Correction style mimicry.} Answers must follow the specific
    rhetorical patterns of official Tunisian corrections, including standard
    French mathematical phrasing (``On a\ldots'', ``Or\ldots'',
    ``Donc\ldots'', ``D'apr\`{e}s\ldots'') and structured proof layouts.

    \item \textbf{Document availability.} No machine-readable corpus of
    Tunisian Bac materials exists. The primary sources are printed booklets,
    handwritten corrections, and scanned pages that must be digitized before
    any retrieval system can operate on them.
\end{enumerate}

\subsection{Approach}
\label{subsec:approach}

We address these challenges through a Retrieval-Augmented Generation (RAG)
pipeline~\cite{lewis2020rag} specifically designed for the Tunisian Bac
mathematics domain. The system comprises:

\begin{itemize}
    \item A \textbf{digitization pipeline} that converts scanned mathematical
    documents into machine-readable LaTeX using multimodal LLMs.

    \item A \textbf{vector database} built on ChromaDB~\cite{chromadb2024}
    with BGE-M3 multilingual embeddings~\cite{chen2024bge_m3}, indexed with
    rich metadata (chapter, year, session, document type).

    \item A \textbf{two-stage retrieval mechanism} that first searches
    corrected exercises for style exemplars, then falls back to course
    material for theorem backing.

    \item A \textbf{prompt compilation layer} that instructs the LLM to mimic
    the correction style when similar exercises are found, or construct
    answers from theorems when they are not.
\end{itemize}

To evaluate the RAG approach, we implement and compare two additional
baselines: a \textbf{prompt-only} system that encodes the entire curriculum
into the system prompt without retrieval, and a \textbf{hybrid} system that
dynamically routes queries between RAG and prompt-only modes based on
retrieval confidence.

\subsection{Contributions}
\label{subsec:contributions}

The main contributions of this thesis are:

\begin{enumerate}
    \item The first publicly available digitized corpus of Tunisian
    Baccalaureate mathematics materials, processed from raw scans into
    structured LaTeX.

    \item A domain-specific RAG architecture with two-stage retrieval
    designed for curriculum fidelity and correction style mimicry.

    \item A systematic comparison of three generation approaches (RAG,
    prompt-only, hybrid) on the Tunisian Bac mathematics domain.

    \item An open-source implementation deployable on Google Cloud Platform,
    including a Streamlit-based interactive tutoring interface.
\end{enumerate}

\subsection{Thesis Structure}
\label{subsec:thesis-structure}

The remainder of this thesis is organized as follows.
Section~\ref{sec:background} reviews the relevant background on large
language models, retrieval-augmented generation, and AI in education.
Section~\ref{sec:methodology} describes the system architecture and
implementation in detail. Section~\ref{sec:experiments} presents the
experimental setup, evaluation methodology, and results.
Section~\ref{sec:discussion} interprets the findings and discusses
limitations. Section~\ref{sec:conclusion} concludes with a summary and
directions for future work.
